% Options for packages loaded elsewhere
\PassOptionsToPackage{unicode}{hyperref}
\PassOptionsToPackage{hyphens}{url}
%
\documentclass[
  10pt,
  a4paper,
]{article}
\usepackage{amsmath,amssymb}
\usepackage{iftex}
\ifPDFTeX
  \usepackage[T1]{fontenc}
  \usepackage[utf8]{inputenc}
  \usepackage{textcomp} % provide euro and other symbols
\else % if luatex or xetex
  \usepackage{unicode-math} % this also loads fontspec
  \defaultfontfeatures{Scale=MatchLowercase}
  \defaultfontfeatures[\rmfamily]{Ligatures=TeX,Scale=1}
\fi
\usepackage{lmodern}
\ifPDFTeX\else
  % xetex/luatex font selection
  \setmainfont[]{DejaVu Serif}
  \setmonofont[]{DejaVu Sans Mono}
  \setmathfont[]{Latin Modern Math}
\fi
% Use upquote if available, for straight quotes in verbatim environments
\IfFileExists{upquote.sty}{\usepackage{upquote}}{}
\IfFileExists{microtype.sty}{% use microtype if available
  \usepackage[]{microtype}
  \UseMicrotypeSet[protrusion]{basicmath} % disable protrusion for tt fonts
}{}
\makeatletter
\@ifundefined{KOMAClassName}{% if non-KOMA class
  \IfFileExists{parskip.sty}{%
    \usepackage{parskip}
  }{% else
    \setlength{\parindent}{0pt}
    \setlength{\parskip}{6pt plus 2pt minus 1pt}}
}{% if KOMA class
  \KOMAoptions{parskip=half}}
\makeatother
\usepackage{xcolor}
\usepackage[margin=20mm]{geometry}
\setlength{\emergencystretch}{3em} % prevent overfull lines
\providecommand{\tightlist}{%
  \setlength{\itemsep}{0pt}\setlength{\parskip}{0pt}}
\setcounter{secnumdepth}{-\maxdimen} % remove section numbering
\usepackage{amsmath,amssymb,amsthm,booktabs,longtable,needspace}
\usepackage{fvextra}
\DefineVerbatimEnvironment{Highlighting}{Verbatim}{commandchars=\\\{\},breaklines=true,breakanywhere=true,fontsize=\small}
\setlength{\emergencystretch}{3em}
\allowdisplaybreaks
\raggedbottom
\widowpenalty=10000\clubpenalty=10000
\providecommand{\Tr}{\operatorname{Tr}}
\providecommand{\tr}{\operatorname{Tr}}
\ifLuaTeX
  \usepackage{selnolig}  % disable illegal ligatures
\fi
\usepackage{bookmark}
\IfFileExists{xurl.sty}{\usepackage{xurl}}{} % add URL line breaks if available
\urlstyle{same}
\hypersetup{
  pdftitle={Separable-reference detector tests},
  pdfauthor={Kirtus Justus},
  hidelinks,
  pdfcreator={LaTeX via pandoc}}

\title{Separable-reference detector tests}
\author{Kirtus Justus}
\date{}

\begin{document}
\maketitle

\subsection{2. Postulates for a new conditional calibration protocol}\label{postulates-for-a-new-conditional-calibration-protocol}

\begin{enumerate}
\def\labelenumi{\arabic{enumi}.}
\tightlist
\item
  \textbf{Known space and reference controls.} The physical system is two qubits with a known tensor decomposition and trusted local Pauli operations. This decomposition and those operations are supplied; they are not derived from a signed-relational substrate.
\item
  \textbf{Stable binary detector.} For each preparation and measurement condition the same unknown Hermitian effect \(E\) acts on the two qubits, with \(0\preceq E\preceq I\). The event complement is \(I-E\). No false-click/missed-click parameterization, diagonal assumption, or classical-only noise assumption is made.
\item
  \textbf{Trusted Bell dephasing.} Immediately before every detector readout, independently randomized local gates implement \[
   \mathcal T(\rho)=\frac14\sum_{g\in\{II,XX,YY,ZZ\}}g\rho g^\dagger.
  \] This map is an average of product unitaries and hence maps separable states to separable states.
\item
  \textbf{Trusted separable reference preparations.} Four source preparations are physically characterized as \(|00\rangle, |++\rangle,|+i,-i\rangle,|+i,+i\rangle\), where \(|\pm i\rangle=(|0\rangle\pm i|1\rangle)/\sqrt2\). They are not merely assigned labels while their Bloch radii are left unknown.
\item
  \textbf{Independent sampling.} For finite-trial guarantees only, trials within each condition are independent with stable outcome laws. The reference and target settings are interleaved or otherwise characterized for drift. Their outcomes are recorded per prepared pair.
\end{enumerate}

\subsection{3. Theorem: separable reference tomography of a general binary effect}\label{theorem-separable-reference-tomography-of-a-general-binary-effect}

Let \(B_j=|b_j\rangle\langle b_j|\) denote the Bell projectors in the order \((\Phi^+,\Phi^-,\Psi^+,\Psi^-)\). The twirl is self-adjoint and projects operators onto the commuting algebra generated by \(XX\) and \(ZZ\). Consequently, \[
 E_{\mathcal T}=\mathcal T^*(E)=\sum_{j=1}^4 e_jB_j,\qquad e_j=\operatorname{Tr}(E B_j)\in[0,1].
\] The event probabilities of the four trusted product reference preparations form \(r=Me\) with \[
 M=\frac12\begin{pmatrix}1&1&0&0\\1&0&1&0\\1&0&0&1\\0&1&1&0\end{pmatrix}.
\] Exact symbolic calculation gives \[
 \det M=\frac18,\qquad
 M^{-1}=\begin{pmatrix}
 1&1&0&-1\\1&-1&0&1\\-1&1&0&1\\-1&-1&2&1
 \end{pmatrix},\qquad \|M^{-1}\|_{\infty}=5.
\] Thus the physically twirled detector response is determined by four separable reference event probabilities. An unrestricted two-qubit Hermitian effect has 16 real parameters, so four settings cannot perform unrestricted effect tomography without the trusted twirl; a full spanning linear tomography would require at least 16 independent input preparations. This is a parameter-count statement about full effect identification, not a universal lower bound for every entanglement witness. \textbf{This is conditional tomography, not a measurement-device-independent claim.} The trusted gate and preparation conditions supply the extra physical information.

\subsection{4. Theorem: exact separable bounds for any twirled binary effect}\label{theorem-exact-separable-bounds-for-any-twirled-binary-effect}

Write the recovered Bell eigenvalues in sorted order \(e_{(1)}\ge e_{(2)}\ge e_{(3)}\ge e_{(4)}\). Then for every separable two-qubit state \(\sigma\), \[
 \boxed{\frac{e_{(3)}+e_{(4)}}2\le \operatorname{Tr}(E\mathcal T(\sigma))\le\frac{e_{(1)}+e_{(2)}}2.}
\] Proof: Local-unitary twirling preserves separability and makes the state Bell diagonal: \(\mathcal T(\sigma)=\sum_jq_jB_j\). By the two-qubit PPT criterion, Bell-diagonal separability is exactly the polytope \(q_j\in[0,1/2],\sum q_j=1\). Linear optimization places \(1/2\) on the two largest coefficients for the maximum, or on the two smallest for the minimum. These bounds are attained by Bell-diagonal separable mixtures of two Bell states of equal weight. Therefore an observed target event probability outside the interval certifies that its original state was entangled under the stated physical controls. Failure to violate the interval is inconclusive: Bell twirling can destroy entanglement.

This is a standard consequence of Bell twirling, PPT separability and linear optimization, not claimed as a new general quantum theorem.

\subsection{5. Negative control A: a physically valid off-diagonal detector fakes a violation without twirling}\label{negative-control-a-a-physically-valid-off-diagonal-detector-fakes-a-violation-without-twirling}

Consider the effect \[
 E=0.1I+0.8B_{\Phi^+}+0.05(|\Phi^+\rangle\langle\Psi^+|+|\Psi^+\rangle\langle\Phi^+|).
\] Its eigenvalues are strictly inside \([0,1]\), so it is a legitimate detector effect. Its Bell diagonal entries are \((0.9,0.1,0.1,0.1)\), whose separable upper envelope \textbf{after twirling} is \(1/2\). Yet the separable state \(|++\rangle\) gives \[
 \operatorname{Tr}(E|++\rangle\langle++|)=\frac{11}{20}=0.55,
\] which would be a false certificate if someone simply treated \(E\) as Bell diagonal. Applying the trusted twirl to the same product state changes the event probability to exactly \(1/2\). The explicit matrix calculations are checked with exact SymPy arithmetic in \texttt{verify\_twirl\_detector.py}.

For the stipulated model-B Bell-diagonal state with \(\Pr(\Phi^+)=653/1280\), the same valid off-diagonal detector gives event probability \(813/1600=0.508125\) after twirling, exceeding its exact separable ceiling by \(13/1600\). This shows that the calibrated approach can detect the specified entanglement without assuming a classical bit-flip model; it does not claim universal detector efficiency.

\subsection{6. Negative control B: trusted twirling does not replace trusted reference-state purity}\label{negative-control-b-trusted-twirling-does-not-replace-trusted-reference-state-purity}

Retain the two original gate-set source radii \(a_A=1/2\) and \(a_B=3/4\), which are unknown to an experimenter using only the native gate-set labels. Let model A's detector effect be \(E_A=B_{\Phi^+}\). Let model B use the legitimate effect \[
 E_B=\frac14\left[I+\frac{XX-YY+ZZ}{(3/2)^2}\right],
\] whose Bell eigenvalues are \((7/12,5/36,5/36,5/36)\). The effects and sources are related by the original similarity gauge, and the Bell twirl commutes with that gauge.

Exact matrix tests verify identical event probabilities for \textbf{all 36 ordered product inputs} constructed from the two model-specific labeled six-axis sources, including after twirling. For the corresponding gate outputs, both predict exactly \(117/320\) clicks, although one output is separable and the other is entangled. Their respective separable upper bounds are \(1/2\) and \(13/36\); the second output exceeds its own bound by \(13/2880\). Thus four reference \emph{labels}, or even all native axis reference labels, do not complete tomography when their physical radii remain uncharacterized.

\textbf{The positive calibration theorem requires trusted physical reference preparations, not just well-conditioned algebra.} This is the same gauge issue seen in gate-set tomography, not a contradiction of the positive theorem.

\subsection{7. Finite-trial theorem and systematic error}\label{finite-trial-theorem-and-systematic-error}

Suppose each of the four reference probabilities and one target probability is estimated from \(M\) independent Bernoulli trials. For an error level \(\alpha\), define \[
 \delta_\alpha(M)=\sqrt{\frac{\log(10/\alpha)}{2M}}.
\] A union bound and Hoeffding imply that all five errors are at most \(\delta_\alpha\) simultaneously with probability at least \(1-\alpha\).

Let \(\widehat e=M^{-1}\widehat r\) and \(b_+(e)=(e_{(1)}+e_{(2)})/2\). Because \(\|M^{-1}\|_\infty=5\) and \(b_+\) is 1-Lipschitz in infinity norm, \[
 \big|b_+(\widehat e)-b_+(e)\big|\le5\delta_\alpha.
\] A valid one-sided size-\(\alpha\) entanglement test is therefore: \[
 \boxed{\text{reject separability only if }\widehat p_{\mathrm{target}}>b_+(\widehat e)+6\delta_\alpha.}
\] For a fixed entangled alternative whose true calibrated effect has positive witness margin \(\Delta=p_{\rm target}-b_+(e)>0\), the same concentration event with failure probability \(\beta\) proves power at least \(1-\beta\) whenever \[
 \Delta>6[\delta_\alpha(M)+\delta_\beta(M)].
\] For the ideal Bell effect \(E=B_{\Phi^+}\), model B has \(\Delta=13/1280\). With \(\alpha=\beta=.05\), \(M=3,698,314\) preparations \textbf{per setting} (four references and one target) satisfies the strict sufficient inequality, for 18,491,570 prepared pairs in total. This is deliberately conservative and \textbf{not an optimized cost}. It exposes the price of calibrating a largely uncharacterized detector to precision at a 1\%-level witness margin. The source's 26,226-pair ideal-detector optimum assumes a different, stronger prior calibration and thus is not directly comparable.

If each of the five settings carries an additional independently bounded systematic probability bias \(\xi\), the same proof yields a valid threshold with \(6(\delta_\alpha+\xi)\). Sufficient power against a nominal alternative needs \(\Delta>6(\delta_\alpha+\delta_\beta+2\xi)\) under worst-case independent biases, or tighter bounds when common-mode drift is characterized. Unbounded setting-dependent drift, state leakage or failures of the twirl invalidate these certificates.

\end{document}
